% !TEX root = chapter-standalone.tex

\section[Isomorphisms]{Isomorphisms}
\label{sec:isomorphisms}



What are identity morphisms good for? One thing we can do with them is define, for any category, the important notion of isomorphism.
This concept describes a way of saying when two objects are ``the same'', even if they are not equal.

\begin{definition}\label{def:isomorphism}
    Let $\CatC$ be a category.
    A morphism $\mora \colon \Obja \mto \Objb$ in $\CatC$ is an \emph{isomorphism} if there exists a morphism $\morb \colon \Objb \mto \Obja$ in $\CatC$ such that
    \begin{equation}
        \mora \mthen \morb = \catid_\Obja
    \end{equation}
    and
    \begin{equation}
        \morb \mthen \mora = \catid_\Objb.
    \end{equation}
\end{definition}

\begin{remark}
    Note that the above definition coincides, for the category $\Set$ of sets and functions, with \cref{def:function-isomorphism}.
    We saw in \cref{ex:bijective-functions-are-isomorphisms} that an isomorphism in the category \Set is the same thing as a bijective function.
\end{remark}

\devel{
    \subsection{Monos and Epis}

    \todotext{J: @J: include this topic at all?}
}

% \linkvideo{spring2021-semicats-cats:summary} % Summary

